\documentclass[10pt,a4paper]{amsart}
\usepackage[margin=19mm]{geometry}
\usepackage{amsmath,amssymb,mathtools}
\usepackage[scr=rsfs,cal=euler]{mathalfa}
\usepackage{xfrac}
\usepackage[unicode,hidelinks,bookmarks=false]{hyperref}
\newcommand{\ie}{i.e.\ }
\newcommand{\cf}{cf.\ }
\newcommand{\resp}{respectively\ }
\newcommand{\Subsection}{\subsection}
\providecommand{\nobreakdash}{\nobreak\hbox{-}}
\setlength{\emergencystretch}{2em}
\multlinegap0pt
\usepackage{amssymb}
\newtheorem{theorem}{Theorem}
\newtheorem{lemma}{Lemma}
\newcommand{\CT}{C_{\mathrm T}}
\newcommand{\CTsimp}{C_{\mathrm T}^{(1)}}
\newcommand{\CTsimpenv}{C_{\mathrm T}^{(1),\leq}}
\newcommand{\ellmax}{\ell_{\max}}
\title{Non-Simple T-Prescriptions Yield T-Complexity Gains Infinitely Often}
\author{Thomas Sch\"urmann}
\date{Source: arXiv:2606.13766v1 (CC BY 4.0)}
\begin{document}
\maketitle

\noindent\emph{Source and licence.} Non-Simple T-Prescriptions Yield T-Complexity Gains Infinitely Often. Version arXiv:2606.13766v1 (CC BY 4.0), distributed by arXiv under the Creative Commons Attribution 4.0 International licence, http://creativecommons.org/licenses/by/4.0/, as recorded in the arXiv licence metadata for this exact version and shown on the abstract page https://arxiv.org/abs/2606.13766v1. Full licence notice: \texttt{licenses/arxiv-2606.13766-CC-BY-4.0.txt}.\par\smallskip

\noindent\textit{Editorial scope.} Counted unit: the article's theorem thm:main (source0.tex lines 106--117). The article's complete original proof of it is delivered with it (source0.tex lines 314--375, one \texttt{proof} environment of 62 lines, which the article places away from the statement). No mathematics is written, repaired or re-derived: the statement and the proof are the article's own lines, in the article's own notation, and no retained line is substituted or re-flowed.  Retained uncounted as the source-local context the proof needs: lines 200--218; lines 258--264.  Nothing was omitted from any retained span, so the package carries no omission record.  No retained line contains a citation, so the wrapper carries no bibliography and no bibliography entry.  No retained line contains a figure environment, an image-inclusion command or drawing code, so no image is delivered and the package declares no asset dependency.  Editorial formatting only, in addition: an AMS \texttt{amsart} preamble that restates the article's own declarations and macros used by the retained lines, namely the environments \texttt{theorem}, \texttt{lemma} on separate counters, together with the standard packages those lines need.  The retained environments print in this document as Theorem 1, Lemma 1 in order of appearance, whereas the article prints the same items with the numbering of its own full document.  The complete native source of the article as distributed in the arXiv e-print for this exact version is bundled unchanged as \texttt{sources/202726/source0.tex}.
\begin{lemma}[Maximum codeword length]
\label{lem:length}
Let
\begin{equation}
P=((p_1,k_1),\ldots,(p_m,k_m))
\end{equation}
be a valid T-prescription. Then
\begin{equation}
\ellmax(P)=1+\sum_{i=1}^{m}k_i|p_i|.
\end{equation}
In particular, if $P$ is simple, then
\begin{equation}
\ellmax(P)=1+\sum_{i=1}^{m}|p_i|
\end{equation}
and
\begin{equation}
\CT(P)=m.
\end{equation}
\end{lemma}

\begin{lemma}[Threshold increments have infinitely many upward jumps]
\label{lem:jumps}
There are infinitely many indices $m$ such that
\begin{equation}
d_{m+1}>d_m.
\end{equation}
\end{lemma}

\begin{theorem}[Infinitely many strict non-simple gains]
\label{thm:main}
For every alphabet size $r\geq 2$, there are infinitely many lengths $N$ such that
\begin{equation}
\CT(N)>\CTsimpenv(N).
\end{equation}
In particular, there are infinitely many lengths $N$ such that
\begin{equation}
\CT(N)>\CTsimp(N).
\end{equation}
Thus non-simple T-prescriptions have strictly larger T-complexity than all simple T-prescriptions at infinitely many finite lengths.
\end{theorem}

\begin{proof}[Proof of Theorem~\ref{thm:main}]
Choose $m$ such that
\begin{equation}
d_{m+1}>d_m.
\end{equation}
Let $P_m$ be a simple prescription with arity $m$ and minimal maximum codeword length:
\begin{equation}
\ellmax(P_m)=L_m.
\end{equation}
Write
\begin{equation}
P_m=((p_1,1),\ldots,(p_m,1)).
\end{equation}
Let $P_m^-$ be the prefix prescription consisting of the first $m-1$ steps. Since $P_m^-$ is simple of arity $m-1$, the definition of $L_{m-1}$ gives
\begin{equation}
\ellmax(P_m^-)\geq L_{m-1}.
\end{equation}
The last simple step from $P_m^-$ to $P_m$ increases the maximum codeword length by $|p_m|$. Therefore
\begin{equation}
|p_m|=\ellmax(P_m)-\ellmax(P_m^-)
       \leq L_m-L_{m-1}=d_m.
\end{equation}
Now keep the first $m-1$ steps and replace only the last simple step $(p_m,1)$ by the non-simple step $(p_m,2)$. This gives the valid prescription
\begin{equation}
\widehat P_m=((p_1,1),\ldots,(p_{m-1},1),(p_m,2)).
\end{equation}
It is valid because $p_m$ is available after the first $m-1$ steps. By Lemma~\ref{lem:length}, the first $m-1$ increments are unchanged, while the last increment is $2|p_m|$ instead of $|p_m|$. Hence its maximum codeword length is
\begin{equation}
N_m=\ellmax(\widehat P_m)=\ellmax(P_m)+|p_m|=L_m+|p_m|.
\end{equation}
Using $|p_m|\leq d_m$ and $d_{m+1}>d_m$, we get
\begin{equation}
N_m\leq L_m+d_m<L_m+d_{m+1}=L_{m+1}.
\end{equation}
Thus no simple prescription with arity $m+1$ or larger can have maximum codeword length at most $N_m$. To see this explicitly, suppose that a simple prescription $Q$ has arity $q\geq m+1$ and $\ellmax(Q)\leq N_m$. Let $Q^{[m+1]}$ be its prefix of length $m+1$. By Lemma~\ref{lem:length}, every later simple step adds a positive word length, so
\begin{equation}
\ellmax(Q^{[m+1]})\leq \ellmax(Q)\leq N_m<L_{m+1}.
\end{equation}
This contradicts the definition of $L_{m+1}$ as the minimal maximum codeword length among all simple prescriptions of arity $m+1$.

Hence every simple prescription counted by $\CTsimpenv(N_m)$ has at most $m$ steps, and so
\begin{equation}
\CTsimpenv(N_m)\leq m.
\end{equation}
On the other hand,
\begin{equation}
\CT(\widehat P_m)=(m-1)+\log_2 3>m.
\end{equation}
Since $\ellmax(\widehat P_m)=N_m$, this gives
\begin{equation}
\CT(N_m)\geq \CT(\widehat P_m)>m\geq \CTsimpenv(N_m).
\end{equation}
There are infinitely many indices $m$ with $d_{m+1}>d_m$ by Lemma~\ref{lem:jumps}. The corresponding lengths $N_m$ are unbounded: since $d_j\geq 1$ for every $j\geq 1$,
\begin{equation}
L_m=L_0+\sum_{j=1}^{m}d_j\geq 1+m\longrightarrow\infty,
\end{equation}
and $N_m\geq L_m$. Hence infinitely many distinct lengths satisfy the claimed strict inequality. Since
\begin{equation}
\CTsimp(N)\leq \CTsimpenv(N),
\end{equation}
the exact-length inequality follows for the same lengths.
\end{proof}

\end{document}
